
Results are presented for each sub-hypothesis in the order of the causal chain: existence (h-e1) $\rightarrow$ expressiveness (h-m1) $\rightarrow$ multi-layer detection (h-m2) $\rightarrow$ failure prevention (h-m3).

\subsubsection{5.1 h-e1: Contract Framework Feasibility}

Three-layer validation proved implementable with minimal overhead. Table 1 summarizes feasibility metrics.

\textbf{Table 1: Contract Framework Feasibility (h-e1)}

\begin{table*}[htbp]
\centering
\resizebox{\dimexpr 2\columnwidth+\columnsep\relax}{!}{%
\begin{tabular}{llllll}
\toprule
Metric & Schema-Only & Contract-Based & Improvement & Threshold & Status \\
\midrule
Detection Rate & 33.3\% (5/15) & 100.0\% (15/15) & +200\% & >50\% & PASS \\
Execution Overhead & 0.01ms & 0.01ms & +0\% & <100ms & PASS \\
Contract Layer LOC & --- & 25 LOC & --- & <50 LOC & PASS \\
False Positive Rate & 0\% & 0\% & +0pp & <10\% & PASS \\
\bottomrule
\end{tabular}
}
\end{table*}

Contract-based validation detected all 15 adversarial violations (100\%) compared to 33.3\% for schema-only, a 200\% improvement exceeding the 50\% threshold. Execution overhead remained negligible (0.01ms per validation), well below the 100ms threshold. Implementation cost was 25 lines of code for the contract layer, within the 50 LOC maintainability bound. Zero false positives confirmed contracts do not over-restrict valid inputs.

\subsubsection{5.2 h-m1: Contracts Force Explicit Checking}

Contract-based validation enforced 100\% of constraints (4/4) while schema-only could not express any (0/4), establishing a 100 percentage point expressiveness gap. Table 2 shows per-constraint coverage.

\textbf{Table 2: Constraint Expressiveness (h-m1)}

\begin{table*}[htbp]
\centering
\resizebox{\dimexpr 2\columnwidth+\columnsep\relax}{!}{%
\begin{tabular}{lllll}
\toprule
Constraint & Type & Schema Coverage & Contract Coverage & Gap \\
\midrule
C1 (No synthetic data) & Keyword pattern & 0\% & 100\% & +100pp \\
C2 (No human eval) & Keyword blacklist & 0\% & 100\% & +100pp \\
C3 (Standard dataset) & Cross-field logic & 0\% & 100\% & +100pp \\
C4 (No new benchmarks) & State-based & 0\% & 100\% & +100pp \\
**Overall** & --- & **0.0\% (0/4)** & **100.0\% (4/4)** & **+100pp** \\
\bottomrule
\end{tabular}
}
\end{table*}

Schema validation checked types (dataset\_name: str) and presence (required fields) but could not restrict string content (C1/C2), cross-field dependencies (C3), or external state membership (C4). Contract-based validation expressed all constraints through field validators (C1/C2), postconditions (C3), and custom validation methods (C4).

The 100 percentage point expressiveness gap (exceeding the 75 percentage point threshold) validates that contracts force explicit checking of constraints schema-only validation leaves implicit.

\subsubsection{5.3 h-m2: Multi-Layer Validation Detects More Violations}

Three-layer validation detected 88\% of violations (22/25) compared to 40\% for schema-only (10/25), a 48 percentage point gap. Table 3 shows layer-wise contributions.

\textbf{Table 3: Detection Rates by Layer (h-m2)}

\begin{table}[htbp]
\centering
\resizebox{\columnwidth}{!}{%
\begin{tabular}{llll}
\toprule
Layer & Violations Detected & Cumulative Rate & Marginal Contribution \\
\midrule
Schema & 10/25 & 40\% & --- (baseline) \\
+ Pattern & +8/25 & 72\% & +32pp \\
+ Contract & +4/25 & 88\% & +16pp \\
\bottomrule
\end{tabular}
}
\end{table}

Layers contributed cumulatively, not redundantly. Schema caught structural errors (10 violations). Pattern layer added 8 semantic violations (keyword detection for C1/C2). Contract layer added 4 compositional violations (cross-field logic for C3/C4). The 48 percentage point gap exceeded the 40 percentage point threshold.

\textbf{Per-constraint detection (Table 4):}

\begin{table}[htbp]
\centering
\resizebox{\columnwidth}{!}{%
\begin{tabular}{llll}
\toprule
Constraint & Violations & Detected & Rate \\
\midrule
C1 (Synthetic) & 5 & 4 & 80\% \\
C2 (Human eval) & 5 & 4 & 80\% \\
C3 (Standard dataset) & 2 & 2 & 100\% \\
C4 (New benchmark) & 3 & 2 & 67\% \\
\bottomrule
\end{tabular}
}
\end{table}

Pattern layer achieved 80\% recall on C1/C2 semantic constraints (met threshold). Contract layer achieved 80\% recall on C3 and 67\% on C4.

\subsubsection{5.4 h-m3: Early Detection Prevents Downstream Failures}

Contract-based validation eliminated 100\% of Phase 4/5 downstream failures (20/20 violations caught at boundaries $\rightarrow$ 0/20 failures) compared to 100\% failure rate for schema-only (0/20 caught $\rightarrow$ 20/20 failures). Table 5 summarizes failure rates.

\textbf{Table 5: Downstream Failure Reduction (h-m3)}

\begin{table*}[htbp]
\centering
\resizebox{\dimexpr 2\columnwidth+\columnsep\relax}{!}{%
\begin{tabular}{lllll}
\toprule
Condition & Boundary Detection & Phase 4/5 Failures & Failure Rate & Reduction \\
\midrule
Schema-only & 0/20 (0\%) & 20/20 & 100\% & --- \\
Contract-based & 20/20 (100\%) & 0/20 & 0\% & **100\%** \\
\bottomrule
\end{tabular}
}
\end{table*}

All 20 injected violations were caught at phase boundaries (Phase $2\rightarrow 3\rightarrow 4$ transitions) under contract-based validation, preventing any violations from reaching Phase 4/5 implementation. Schema-only validation missed all 20 violations, resulting in 100\% Phase 4/5 failure rate as violations propagated undetected until implementation discovered infeasibility.

\textbf{Per-constraint failure analysis (Table 6):}

\begin{table*}[htbp]
\centering
\resizebox{\dimexpr 2\columnwidth+\columnsep\relax}{!}{%
\begin{tabular}{llllll}
\toprule
Constraint & Injected & Boundary Caught & Schema Failures & Contract Failures & Reduction \\
\midrule
C1 & 5 & 5/5 (100\%) & 5/5 (100\%) & 0/5 (0\%) & 100\% \\
C2 & 5 & 5/5 (100\%) & 5/5 (100\%) & 0/5 (0\%) & 100\% \\
C3 & 5 & 5/5 (100\%) & 5/5 (100\%) & 0/5 (0\%) & 100\% \\
C4 & 5 & 5/5 (100\%) & 5/5 (100\%) & 0/5 (0\%) & 100\% \\
\bottomrule
\end{tabular}
}
\end{table*}

The 100\% failure reduction exceeded the 80\% threshold. Zero false positives on 80 valid hypotheses confirmed contracts do not over-restrict feasible research.

\subsubsection{5.5 Summary of Key Results}

\begin{enumerate}
\item \textbf{Feasibility (h-e1):} Three-layer validation implementable with 25 LOC, 0.01ms overhead, 200\% detection improvement.
\item \textbf{Expressiveness (h-m1):} 100 percentage point gap---contracts check 4/4 constraints, schema 0/4.
\item \textbf{Multi-layer detection (h-m2):} 48 percentage point gap---88\% versus 40\%, with cumulative layer contributions (schema 10, pattern +8, contract +4).
\item \textbf{Failure prevention (h-m3):} 100\% reduction---all violations caught at boundaries, zero Phase 4/5 failures.
\end{enumerate}

All sub-hypotheses met their success criteria. Primary prediction ($\geq 80$\% failure reduction) validated at 100\%. Secondary prediction ($\geq 95$\% detection) achieved partial support (88\% via 40 percentage point gap criterion).

The causal mechanism holds: contract specification $\rightarrow$ forced explicit checking (100 percentage point gap) $\rightarrow$ multi-layer detection (88\% versus 40\%) $\rightarrow$ prevented failures (100\% reduction).
